\section{Tor Hidden Service}

Konfiguration in \texttt{/etc/tor/torrc}:

\begin{lstlisting}[language=bash]
# torrc Konfiguration
HiddenServiceDir /var/lib/tor/hidden_service/
HiddenServicePort 80 127.0.0.1:80

# Service starten
systemctl restart tor

# .onion-Adresse auslesen
cat /var/lib/tor/hidden_service/hostname
\end{lstlisting}

\section{Tor-Artefakte suchen}

Im forensischen Image nach Tor-Spuren suchen:

\begin{lstlisting}[language=bash]
# Tor-Dateien im Dateisystem
fls -pr /dev/loopXpN | grep -i 'tor'

# Prefetch: Wurde Tor ausgefuehrt?
find /mnt -type f -iname "*.pf" | grep -i tor
sccainfo /mnt/Windows/Prefetch/TOR.EXE-*.pf

# Onion-URLs in Browser-History (Firefox-basiert)
fls -pr /dev/loopXpN | grep -i 'places.sqlite' | grep -i tor
icat /dev/loopXpN <INODE> > /tmp/tor_places.sqlite
sqlite3 /tmp/tor_places.sqlite \
  "SELECT url FROM moz_places WHERE url LIKE '%.onion%';"

# Strings-Suche nach .onion-Adressen
strings /dev/loopXpN | grep '\.onion'
grep -r "\.onion" /mnt/Users/ 2>/dev/null
\end{lstlisting}

\section{SquashFS-Container}

Versteckte Daten in komprimierten Containern am Ende eines Images:

\begin{lstlisting}[language=bash]
# Container finden
binwalk image.dd
# -> Squashfs filesystem at offset X

# Container extrahieren
dd if=image.dd bs=1 skip=OFFSET count=SIZE \
  of=/tmp/container.squashfs

# Entpacken
unsquashfs /tmp/container.squashfs
ls squashfs-root/
\end{lstlisting}

Beispiel: mkp224o (Tor Vanity Address Generator) — generiert \texttt{.onion}-Adressen mit bestimmten Präfixen.
Git-Repos im Container können Aktivitäten dokumentieren.

\section{Versteckte ZIP-Archive}

In dd-Images eingebettete Archive finden:

\begin{lstlisting}[language=bash]
binwalk image.dd
# -> Zip archive at offset X, size Y

dd if=image.dd skip=OFFSET bs=1 count=SIZE \
  of=/tmp/hidden.zip
unzip /tmp/hidden.zip
\end{lstlisting}
